\section{Related Works}\label{s:related_work}

\noindent\paragraph{Model Lineage and Independence Testing.}
The growing population of open-source LLMs has spurred research into determining LLM lineage—whether two models are independently trained or one is trained from another —which is crucial for intellectual property and model safety. Recent works use distinct approaches: output-based methods use output comparison~\citep{Nikolic2025model}, fingerprint queries~\citep{jin2024proflingo}, and lexical cues~\citep{sun2025idiosyncrasies, wu-etal-2023-llmdet} to compare models without requiring parameter access; watermarking techniques inject
provenance information during training or decoding~\citep{wang2019attacks, Uchida2017embedding, Gu2022WatermarkingPL, jiang2025stealthink}; whereas white-box methods perform hypothesis testing~\citep{zhu2025independence} or fingerprinting using metrics derived from model weight structures~\citep{zeng2024HuRef, zeng2025awm, wang2025ghost, yoon2025intrinsic}. 

% Furthermore, recent work~\citep{tong2025seedprints} also suggests that random initialization biases can be "biometric" identifiers throughout the model lifecycle. These methods have established a rigorous framework for understanding intrinsic structures of LLMs and their lineages.


% %
% There are several distinct approaches for model relation inference. Output-based methods require no parameter access. They compare models by issuing random queries to test statistical output similarity~\cite{Nikolic2025model}, using adversarial-style fingerprint queries~\citep{jin2024proflingo}, or training classifiers that exploit stylistic and lexical cues in generated text~\cite{sun2025idiosyncrasies, wu-etal-2023-llmdet}. ~\cite{gao2024model} formalizes when two models should be regarded as computing the same function by characterizing a wide class of output-preserving reparameterizations and then testing functional equality under these symmetries.
% Watermarking techniques instead inject provenance information during training or decoding. Weight-level watermarks embed cryptographic or structured signals directly into model parameters~\cite{wang2019attacks, Uchida2017embedding}, whereas text-level watermarks modify sampling to imprint detectable statistical patterns into outputs~\cite{Gu2022WatermarkingPL, jiang2025stealthink}.
% %
% Another category of works consider model weights.  ~\cite{zhu2025independence} instead treats provenance as a statistical hypothesis test on the model weights. Assuming permutation-equivariant training dynamics, they construct test statistics that yield p-values for the null hypothesis that two models were trained independently. The most relevant work to us is ~\cite{zeng2024HuRef}, which proposes a weight-based fingerprinting approach. By deriving permutation-invariant terms from attention and feed-forward matrices and stacking them into a fixed-size tensor, they map the tensor through a learned encoder and pretrained StyleGAN2 generator to produce a human-readable fingerprint image. Cosine similarity between these invariant tensors (ICS) reliably identifies the base model among dozens of offspring LLMs.


\paragraph{Model zoos and LLM embeddings.} 
There has been growing interest in exploring representation-based approaches to capture model dependence, properties, and relationships, enabling efficient organization and management of large model zoos. Relevant directions include LLM embeddings~\citep{zhuang2025embedllm}, weight-space learning~\cite{horwitz2025Atlas, horwitz2025TreeHeritage, gueta2023KnowledgeRegion}, and loss-landscape-based grouping~\cite{gueta2023KnowledgeRegion, schurholt2025ModelZoo}. LLM embedding methods~\citep{zhuang2025embedllm} learn compact and model-level representations that can be used to predict downstream performance across tasks without additional inference. In contrast, weight-space learning approaches operate directly on model parameters. For example, Horwitz \emph{et al.}~\cite{horwitz2025Atlas} combine weight-space distances with model metadata to construct an ``atlas’’ of open models, enabling similarity search and provenance analysis within large repositories. A complementary line of work~\citep{schurholt2025ModelZoo} considers loss landscape properties and phase-based information~\citep{Yang2021TaxonomizingLossLandscapes}, such as mode connectivity and CKA similarity, to partition models into distinct regimes and assess model similarity. 
% Overall, these studies demonstrate multiple pathways for LLM representations and efficiently structuring large collections of models.


%Alternative approaches have been introduced to capture model dependence, properties, and relationships. Relevant approaches include representation learning~\citep{zhuang2025embedllm}, weight-space learning~\cite{horwitz2025Atlas, horwitz2025TreeHeritage, gueta2023KnowledgeRegion} and grouping by loss landscapes ~\cite{gueta2023KnowledgeRegion, schurholt2025ModelZoo}. 
% %
% Horwitz \emph{et al.}~\cite{horwitz2025Atlas} draws on both weight-space learning and available model metadata to chart an ``atlas'' of open models. Model dependence is determined by calculating the euclidean distance between 100 randomly chosen, corresponding, neurons from a given pair of models. Model metadata like HuggingFace upload timestamps and tensor types are used to extract more specific information about model relationships and provenance. To find models similar to an unlabelled model, Horwitz \emph{et. al.}~\cite{horwitz2025Atlas} propose implementing a simple KNN algorithm with labelled models in their existing atlas. Several assumptions were made in creating this atlas, such as permutation equivariant transformations between parent and child models, and specific HuggingFace user behavior when uploading models and checkpoints. 
% %
% Schürholt \emph{et al.}~\cite{schurholt2025ModelZoo} propose a method rooted in statistical physics, where model similarity is related to its loss landscape metrics and its “phase” in temperature vs. relative-load space, as defined by Yang \emph{et al.}~\cite{Yang2021TaxonomizingLossLandscapes}. Specifically, the phases are defined by the mode connectivity and CKA similarity between pairs of models. Mode connectivity and CKA partition models into five main phases: underfitted and undersized, underfitted, undersized, overfitted, and generalizing. Models within a defined phase are very similar and those in different phases exhibit qualitatively different properties and performances. The boundaries between phases are distinct, with abrupt changes in properties between phases. It is worth noting that while Schürholt et al.~\cite{schurholt2025ModelZoo} uses these two metrics to define model phases, there are other ways to measure loss landscape properties~\cite{andriushchenko2023sharpness, kwon2021asam}. Nevertheless, this study effectively demonstrates a way to use weight matrix analysis to partition models by their properties.


\noindent\paragraph{HT-SR Theory.} HT-SR theory is a theoretical framework based on random matrix theory for analyzing weight matrices geometry of neural networks~\citep{mahoney2019traditional, martin2021implicit, couillet2022random, liao2025random}. HT-SR theory was proposed
based on the observation that well-trained, state-of-the-art DNNs often exhibit Heavy-Tailed structures in the Empirical Spectral Density (ESD)
of each layer. In the meantime, several rigorous theories in SGD relating HT phenomena to generalization
performance were established, providing further theoretical support for HT-SR theory~\citep{gurbuzbalaban2021heavy, hodgkinson2021multiplicative, hodgkinson2022generalization, simsekli2019tail, simsekli2020hausdorff}. Recent works further extend theoretical foundations, and show that HT-SR Theory can be used to assess the training quality of modern DNN models across various domains (CV, NLP, etc.) without accessing any training or test data~\citep{yang2023test, martin2021implicit}. Based on this insight, recent works have further proposed HT-SR theory-based methods for practical neural network applications, including hyperparameter tuning~\citep{zhou2023temperature, liu-etal-2024-model, hu2025eigenspectrum}, model pruning~\citep{lu2024alphapruning,hu2025eigenspectrum}, and more. In this work, we employ HT-SR theory in analyzing the geometric similarities of weight matrices of LLMs.


